\scanpage{31}

l'aide de la notion de \uline{profondeur}.

Les résultats \inserted{ainsi} obtenus sont l'outil technique
principal pour un \deleted{\unsure{développ.}}
\inserted{développement} systématique de
\deleted{propriétés} \inserted{certains théorèmes}
du type de ``Lefschetz'', comparant les propriétés géométriques
d'un schéma $X$ à celles de certains \inserted{sous-schémas}
diviseurs $Y$, \deleted{(section hyperplane)}
\deleted{\unsure{de} $X$} \inserted{\unsure{sur} $X$},
\uline{via} les propriétés du complété formel $X_{/Y}$, \unsure{etc}.
\deleted{Ils \unsure{présentent} \unsure{un} \illegible{2}
\unsure{l'observation} \unsure{faite} \unsure{plus} \unsure{explicite}.}
Ils sont applicables par exemple
\inserted{(après des \unsure{procédés usuels} de passage aux
\unsure{cônes projetants})}
au cas d'une variété projective $X$ et une section hyperplane
convenable, moyennant des hypothèses \deleted{de \unsure{rang}}
convenables sur la profondeur de $X$.
\deleted{\unsure{Ces} \unsure{résultats}}

\inserted{Dans ce cas}
\deleted{\unsure{nous aurons donc} à préciser.}
\uline{Le théorème} \deleted{\unsure{est dû}}
par H. Grauert qui avait d'ailleurs
\deleted{\unsure{été} \unsure{démontré} par \unsure{nous}}
(\deleted{\unsure{Bourbaki}} dès 1958,
(utilisant les th. de \unsure{nullité} de Serre [\ ]) [\ ])
\deleted{\unsure{et montrant que dans la situation précédente}}
\deleted{\unsure{il} \unsure{était} \illegible{2}}

(i) \inserted{Pour tt $F$ loc. libre,}
\[
H^i(X,F)\longrightarrow H^i(\hat X,\hat F)
\quad\text{est bijectif}
\]
pour des \deleted{petits} $i\leq n$
\deleted{[$n$ \unsure{dépendant de}]}
\inserted{$\unsure{\leq}\operatorname{prof}F-2$}.

(ii) si un faisceau formel loc. libre
$\unsure{\mathfrak F|\mathcal F}$ sur $X_{/Y}$,
\deleted{$\unsure{\hat{\mathfrak F}}$ \unsure{est}}
\inserted{et si $\operatorname{prof}(\unsure{O_{X}})\geq2$,}
alors pour $n$ \deleted{\unsure{corps de sections}}
grand, $\unsure{\mathfrak F|\mathcal F}(n)$ est engendré
par ses sections [\unsure{et}] sur un ens. fini --
\deleted{[et \unsure{pour les sections}, \unsure{voir}
\unsure{N°} \unsure{d'un} \unsure{Grothendieck},
\unsure{cela} \unsure{est} \unsure{utile}]}
comme nous montrerons],
(\unsure{Cela qui} implique \inserted{p. ex.}
déjà ce th de Lefschetz pour le groupe fondamental annoncé dans [\ ]).
$\unsure{\rho|p}$

\margin{$\unsure{\rho|p}$ La théorie développée ici et
\unsure{plus loin} en précisera et en généralisera les résultats
de Grauert.}

\margin{30}
